%-*-latex-*-
\input{myassignmentpreamble}
\input{ciss358}
\input{yliow}
\renewcommand\TITLE{Assignment 5}

\begin{document}
\topmatter

\textsc{Objectives}
\begin{itemize}
  \li Solve problems using divide and conquer technique.
  \li Prove results on asymptotics.
  \li Compute asymptotics without using Master Theorem.
\end{itemize}

For programming problems such as Q1, put your code in
\verb!a05q01/!.
Some questions require writing proofs/computations.
Enter then in \verb!q03.tex!, etc.

%============================================================================== 
\newpage
Q1.
Implement the closest pair of points algorithm.

Here's a test case to fix output
\begin{console}[commandchars=\\\{\}]
\userinput{5 2 3 2 10 10 10 11 10 11 11}
1 2 10 10 11 10 11 10 11 11
\end{console}
Explanation of inputs/outputs:
\begin{enumerate}
  \li \verb!5 2 3 2 10 10 10 11 10 11 11!:
  There are \verb!5! points.
  The points are $\verb!(2, 3)!$, $\verb!(2, 10)!$, etc.
  \li \verb!1 2 10 10 11 10 11 10 11 11!:
  The minimal distance is \verb!1!.
  \verb!2! pairs of points are found.
  The pairs of points are
  $\verb!(10, 10)!-\verb!(11, 10)!$ and
  $\verb!(11, 10)!-\verb!(11, 11)!$.
  Sort the points by dictionary order.
\end{enumerate}

Your points should have \verb!double! coordinates.
For computational geometry, it's common to have a \verb!struct! for
points:
\begin{console}
struct Point
{
    double x, y;
};
\end{console}

%============================================================================== 
\newpage
Q2.
You’re consulting for an investment company.
They have the following type of problem that they want to solve over and over.
A typical instance of the problem is the following.
They’re doing a simulation in which they look at $n$ consecutive days of a
given stock, at some point in the past.
The days are numbered $i = 0, 1, ..., n - 1$.
For each day $i$, they have a price $p[i]$ per share for the stock on that day.
Suppose during this time period, they wanted to buy shares on some day
and sell all these shares on some later day.
They want to know when should they have bought and when
should they have sold in order to have made as much money as possible.
If there was no way to make money during the $n$ days, your program
should report it as well.

For example, suppose $n = 3$, $p[0] = 9$, $p[1] = 1$, $p[2] = 5$.
Then your program should print \verb!1 2!, i.e., \lq\lq buy on day 1,
sell on day 2".
Buying on day 1 and selling on day 2 means they
would have made \$4 per share, the maximum possible for that period.
Here's the test case to fix input/output:
\begin{console}[commandchars=\\\{\}]
\userinput{3 9 1 5}
1 2
\end{console}
For the cases when they cannot make any profit, here's a test:  
\begin{console}[commandchars=\\\{\}]
\userinput{3 3 2 1}
None
\end{console}
and another:
\begin{console}[commandchars=\\\{\}]
\userinput{3 1 1 1}
None
\end{console}


%============================================================================== 
\newpage
Q3.
Let $T(n)$ satisfy the following recurrence:
\[
T(n)
= 7 T(n/2) + An^2 
\]
\begin{enumerate}[label=(\alph*)]
\item
  Let $n = 2^k$. Derive $T(n)$ as a sum of power series expressions
  of $n$.
  The expresssion should be of the form
  \[
  T(n) = [?] T(1) + [?] A
  \]
  The $[?]$ should be expressions in $n$ (not $k$).
  Whenever possible, write the expression as $n^?$.
  For instance it's OK to write
  \[
  T(n) = n^{\log_{100} {200}} T(1) + \frac{n^{\log_{111} 222} - 4}{3} A
  \]
\item
  Write down a claim of the form $T(n) = O(g)$ where $g$ is deduced from (b).
  Your $g$ should be of the form $n^\alpha$ or $n^\alpha \log^\beta n$. 
\item
  Prove your claim $T(n) = O(g)$ (from (b)) using induction.
  Make sure you state your $N$ and $C$.
\end{enumerate}

For this question, refer to section \textsc{121.3 Exponentiation} from
the notes on divide and conquer for examples
on how to compute the power series to guess the big-O
and examples on how to then prove the big-O is correct using induction.
There are more computations that involves substitutions in order to guess
the runtime in the chapter Asymptotics
\textsc{123.10 Master theorem}. 

(Obviously you should \textit{not} use the Master Theorem.)

For (b), you have to pick of the two function [?] in
\[
T(n) = [?] T(1) + [?] A
\]
grows the fastest to be your $g(n)$.
The [?] next to $A$ is usually simplified using the geometric sum
formula. So it would be something like
\[
\frac{u^{v} - 1}{u - 1}
\]
Of course you can use
\[
\frac{u^v}{u - 1}
\]
as an upper found for this function and choose between this and
$[?] T(1)$.
For sure you want to write both of them as expressions involving $n$.
Once you have identified which of the two is larger, you will get
an plausible upper bound for $T(n)$.
For instance say for $n = 2^k$, you finally get
\[
T(n) = n^{2.2} T(1) + \frac{1}{10} n^{3.3} A
\]
Then you suspect
\begin{align*}
  T(n)
  &= n^{2.2} T(1) + \frac{1}{10} n^{3.3} A \\
  &\leq n^{3.3} T(1) + \frac{1}{10} n^{3.3} A \\
  &\leq \left( T(1) + \frac{1}{10} A \right) n^{3.3} \\
\end{align*}
for all $n$, and not just $n = 2^k$.
Of course the next step is then to prove
$T(n)$ is bounded above by $Cn^{3.3}$ asymptotically for some $C$.
Note the next VERY important point:
You need not choose $C$ to be $\left( T(1) + \frac{1}{10} A \right)$.
In fact usually this won't work.
Try it out on your own in a proof by induction using this as your $C$
and you will see why it will most likely not work.
Instead, you want to try using
\begin{align*}
  T(n)
  &\leq \left( T(1) + [?] A \right) n^{3.3} \\
\end{align*}
where $[?]$ is \textit{larger} than $1/10$.
For instance, maybe 
\begin{align*}
  T(n)
  &\leq \left( T(1) + 2 A \right) n^{3.3} \\
\end{align*}


\SOLUTION


(a)
$T(n) = 7T(\frac{n}{2}) + An^2$\\

$T(2^k) = 7T(2^{k-1}) + A2^{2k}$\\

$T(2^k) = 7^kT(1) + A2^{2k}(\frac{7^0}{(2^2)^0} + \frac{7^1}{(2^2)^1} + \cdots + \frac{7^k}{(2^2)^k})$\\

$T(2^k) = 7^kT(1) + A2^{2k}(\sum^{k}_{i=0}\frac{7^i}{(2^2)^i})$\\

$T(2^k) = 7^kT(1) + A2^{2k}(\sum^{k}_{i=0}7^i \sum^{k}_{i=0}(2^2)^i)$\\

$T(2^k) = 7^kT(1) + A2^{2k}(\frac{1-7^k}{-6})$\\

$T(2^k) = 7^kT(1) + A2^{2k}(7^k)$\\

$T(2^k) = 7^{\lg n}T(1) + An^2(\frac{1-\frac{7^{\lg n}}{n^2}}{1-\frac{7}{4}})$\\


(b) $T(n) = O(n^2\cdot 7^{\lg n})$

(c)


%============================================================================== 
\newpage
Q4.
Prove the following.
In the following $f, f_i, g, g_i, h$ are functions of $n \in \N$.
You can assume they are nonzero. 
Assume $c, c_i$ are constants, which you can assume are nonzero.
\begin{enumerate}[label=(\alph*)]
\item
  $f = O(g), g = O(h) \implies f = O(h)$
\item
  $f_1 = O(g_1), f_2 = O(g_2)
  \implies f_1 + f_2 = O ( \max(|g_1|, |g_2|) )$
\item
  $f_1 = O(g), f_2 = O(g) \implies f_1 + f_2 = O (g)$
\item
  $f = O(cg)$ $\iff$ $f = O(g)$
\item
  $cf = O(g)$ $\iff$ $f = O(g)$
\end{enumerate}

\SOLUTION

%-*-latex-*-

(a)
Let $f = O(g)$ and $g = O(h)$.
Therefore there are $N_1, C_1, N_2, C_2$ such that
if $n \geq N_1$,
\[
|f(n)| \leq C_1 |g(n)|
\]
and if $n \geq N_2$,
\[
|g(n)| \leq C_2 |h(n)|
\]
Let $N_3 = ?$ and $C_3 = ?$.
Then for $n \geq N_3$,
[...]

(b)
[I proved this in class. Check your notes and write up the proof
  properly.]


(c)

(d)
($\implies$)


($\impliedby$)

(e)
($\implies$)


($\impliedby$)




%===============================================================================
\newpage
Q5. 
Let $p$ and $q$ be polynomials of degree $d$ and $e$ respectively.
Prove the following.
\begin{enumerate}
\item[(a)] $p \sim q$
  iff $d = e$ and leading coefficient of $p$ and $q$ are the same.
\item[(b)] $p = \Theta(n^d)$
\end{enumerate}
You can assume facts about $\sim$ 
(Proposition 123.3.6 -- Proposition 123.3.8).

\SOLUTION

%-*-latex-*-

Let
$p(n) = p_dx^d + p_{d-1}x^{d-1} + \cdots + p_0x^0$
and
$q(n) = q_ex^e + q_{e-1}x^{e-1} + \cdots + q_0x^0$.

(a)
$(\implies)$
Let $p \sim q$.
[...]
Hence $d = e$ and $p_d = q_e$.

$(\impliedby)$
Assume $d = e$ and $p_d = q_e$.
[...]
Hence $p \sim q$.

(b)


\end{document}
